\part{Unestablished stronger claims and source boundaries}
\section{What the checked results do not establish}

This part contains no finding presented as a theorem.  It records the
precise propositions that the task's suggestions could otherwise be read
as asserting, and the limits of the supplied evidence.  A counterexample
to a formal inference is not a counterexample to a geometric conjecture.

\begin{longtable}{@{}p{0.12\textwidth}p{0.35\textwidth}p{0.46\textwidth}@{}}
\toprule
ID & Unestablished assertion & Exact boundary or required additional datum\\
\midrule
\endfirsthead
\toprule ID & Unestablished assertion & Exact boundary or required additional datum\\
\midrule\endhead
U-01 & General Hodge standard positivity, or algebraic Lefschetz correspondences, from sign-only purity. & HS-05 and SG-01--03 refute the proposed formal implication. A construction of the specified positive form and its algebraic adjoint remains additional geometry. BRIDGES/Q05 is unchanged.\\
U-02 & Every algebraic class on an arbitrary CM abelian variety is a Lefschetz class. & SG-04 proves this only for the specified CM elliptic powers and their good reductions. Milne's conditional general theorem retains its actual Hodge-conjecture hypothesis.\\
U-03 & A higher arithmetic Hodge-index theorem, or a corrected general cote-44 height identity, from steps 5 and 13. & AR-01 refutes the displayed raw identity; AR-02--03 refute the proposed sign inference. A corrected theorem must specify admissible cycle extensions, vertical corrections, metrics when arithmetic, and comparison with the canonical self-duality.\\
U-04 & A normalized, lattice-preserving identification of the intended automorphic lift chain with a Lefschetz realization. & SL-05 and SL-07 rule out the automatic lattice-free and untwisted identifications. A new comparison would require a specified realization, Tate normalization, integral intertwiner, and equality of the actual operators and lattices. No such comparison is asserted.\\
U-05 & Exact valuations of a particular nonzero iterated theta period determined solely by the lift count. & SL-06 gives the actual conditional local bounds and explicit numerical contributions. Its inputs include highest weights, primes, denominator bounds and entry estimates. Wilson computes an entirely specified exterior map; the example does not select a nonzero period.\\
U-06 & The original full SA06 C.1f--q proofs, the degree-14 Lawson nonvanishing in E.3u, or the original handwriting in the struck remarks. & Only the supplied intake statements and transcriptions were available. They are not independently upgraded. The present purity assertions do not depend on those absent proofs; the proposed textual repairs require the images.\\
U-07 & A closed Lean certificate for the reconstructed generic-path proof. & The archive was inspected at declaration and selected proof-body level, not built. Source definitions and implication theorems are distinguished from supplied premises. All theorem proofs used here are mathematical proofs with their explicit foundational inputs.\\
\bottomrule
\end{longtable}

The distinction in U-04--05 is part of the answer to the comparison
question: the stated unqualified identification is false, and the two
actual numerical statements have been computed with their separate
parameters.  An additional normalized integral comparison has not been
invented, assumed, or disproved merely because it is not supplied.

\section{Dependencies, verification, and programme consequences}

The main implication chain is
\[
 \begin{gathered}
 \text{Artin vanishing + generic weak Lefschetz + duality}\\
 \Downarrow\\
 \text{SG-08: an effective Gysin kernel and a tensor-power summand}\\
 \text{Noether normalization + HS-03 (Parseval) + HS-01--02}\\
 \Downarrow\\
 \text{SG-09: affine middle upper bound}\\
 \Downarrow\\
 \text{SG-11: compact-support upper bounds}\\
 \Downarrow\\
 \text{SG-10: smooth proper purity by duality.}
 \end{gathered}
\]
There is no arrow here from a trace multiplicity to a positive
intersection form.  Cote 16's evanescent machinery helps locate the
effective kernel; its index and cycle-conjecture assertions are separate.
The only use of the stipulated global RH in these findings is the
change of variable for the explicitly identified Tate $L$-factor in
SG-02.  No general positive-weight number-field $L$-function is converted
into a Dirichlet $L$-function by this argument.

The original BRIDGES/Q03 row is replaced by a literal refutation and
proved precise corrections (Q3-01, Q3-04--07).  BRIDGES/Q04 remains the
known weight-to-coinvariants consequence.  BRIDGES/Q05 remains open in
its general form.  The present SA06 C.1 conclusions are realization
statements: they leave all independent cycle-class-kernel issues in
their proper scope.  H8's exterior-algebra content is corrected and
computed in SL-01--04 and SL-08.  H9's displayed uncorrected higher-cycle
identity is refuted by AR-01.  No statement becomes conditional on
Masullo here, since no specific Masullo theorem is an input.

Every complete proof in Part I received a full reading after its
construction; the principal sign-only theorem and the main counterexamples
also received independent reconstruction by another agent.  The exact
events, including corrected drafts, appear in \texttt{CHECK\_EVENTS.tsv}.
All agents used the inherited GPT-6 Astra Pro model; the runtime did not
expose a separate effort setting, and the log does not manufacture one.
The integer matrix checks are reproducible with the included Python
script and tables.  They are computations, not Lean certificates.
Literature statements inspected but not re-proved are logged with that
weaker event type.  The dated ledger records actual gate changes;
preliminary notes are preserved as working history and explicitly
superseded where the completed geometric proof closed a former gap.

\begin{thebibliography}{99}
\bibitem{Alpoge}
L.~Alp\"oge, \emph{Remark}, supplied four-page note dated 20 November 2026.
Step 5: p.~3, footnote 3 (timeline steps 3 and 5); step 10: divisibility
of the iterated lift; step 13: slope inequality.
\url{https://alpo.ge/remark.pdf}. Supplied wrappers
\texttt{ALPOGE\_Main.lean}, \texttt{ALPOGE\_Statements.lean}; retrieved
source/declaration locators are given in the body and evidence notes.

\bibitem{Archive}
A.~Grothendieck, supplied archive transcriptions: cote 16, scans 014--018,
034--040, 049--050; cote 44, scans 001--009; cote 34, scans 053--064 and
074. The exact file inventory and hashes are in \texttt{SOURCE\_MANIFEST.tsv}.
Original scan images were not included.

\bibitem{SGA1}
A.~Grothendieck et al., \emph{Rev\^etements \'etales et groupe fondamental}
(SGA 1), LNM 224, Springer, 1971; corrected re-edition,
\url{https://arxiv.org/pdf/math/0206203}.
Expos\'e V, Th.~4.1 and Prop.~6.9; Expos\'e IX, Th.~6.1.

\bibitem{SGA4}
M.~Artin, A.~Grothendieck, J.-L.~Verdier et al., \emph{Th\'eorie des topos
et cohomologie \'etale des sch\'emas} (SGA 4), LNM 269, 270, 305.
Expos\'e XI, 3.1, 3.3; XIV, Th.~3.1; XVII, 5.4.3; XVIII (duality).
The uses and source-check depth are recorded in the check table.

\bibitem{SGA4half}
P.~Deligne et al., \emph{Cohomologie \'etale} (SGA $4\frac12$), LNM 569,
Springer, 1977, \emph{Th\'eor\`emes de finitude}, 1.9 and 4.3.

\bibitem{BBD}
A.~Beilinson, J.~Bernstein, P.~Deligne, \emph{Faisceaux pervers},
Ast\'erisque 100 (1982), 4.1.1--4.1.3 and 4.2.8.
\url{https://publications.ias.edu/sites/default/files/Faisceaux%20pervers.pdf}.

\bibitem{WeilI}
P.~Deligne, \emph{La conjecture de Weil. I}, Publ. Math. IHES 43 (1974),
273--307. (1.5.1), (1.5.4), Th.~1.6, (1.7), (2.3), (2.14),
\S\S4.1, 5.9.
\url{https://www.numdam.org/article/PMIHES_1974__43__273_0.pdf}.

\bibitem{WeilII}
P.~Deligne, \emph{La conjecture de Weil. II}, Publ. Math. IHES 52 (1980),
137--252, Th.~3.3.1 and Cor.~3.3.5.
\url{https://publications.ias.edu/sites/default/files/Number40.pdf}.

\bibitem{KatzACT}
N.~M.~Katz, \emph{Affine Cohomological Transforms, Perversity, and
Monodromy}, J. Amer. Math. Soc. 6 (1993), 149--222.
Th.~3.4, Cor.~3.4.1, Cor.~3.4.3 and Lemma 3.5; Appendix A.3--A.5.
\url{https://web.math.princeton.edu/~nmk/old/ACT.pdf}.

\bibitem{StacksNN}
The Stacks Project, \emph{Noether normalization}, Lemma 10.115.4,
Tag 00OY, \url{https://stacks.math.columbia.edu/tag/00OY}.

\bibitem{KatzLang}
N.~M.~Katz and S.~Lang, \emph{Finiteness theorems in geometric classfield
theory}, Enseign. Math. 27 (1981), 285--319, Th.~1, proof in \S II,
Ths.~1(bis), 1(ter).
\url{https://web.math.princeton.edu/~nmk/old/AbsFinThm.pdf}.

\bibitem{MilnePolarizations}
J.~S.~Milne, \emph{Polarizations and Grothendieck's Standard Conjectures},
Ann. of Math. 155 (2002), 599--610. Published-version preprint,
Prop.~3.2, Th.~3.3, Remarks 3.7--3.8.
\url{https://www.jmilne.org/math/articles/2002aP.pdf}.

\bibitem{MilneAV}
J.~S.~Milne, \emph{Abelian Varieties}, version 1.10 (2008).
Th.~2.3; Th.~11.1(b); Prop.~13.1; Th.~13.2; Th.~14.1 and
Lemmas 14.2--14.3.
\url{https://www.jmilne.org/math/CourseNotes/AV110.pdf}.

\bibitem{SutherlandEC}
A.~V.~Sutherland, \emph{18.783 Elliptic Curves}, MIT OpenCourseWare,
Fall 2025, Lecture 13, Ths.~13.3, 13.18 and Cor.~13.20.
\url{https://ocw.mit.edu/courses/18-783-elliptic-curves-fall-2025/mit18_783_f25_lec_full.pdf}.

\bibitem{MilneJacobian}
J.~S.~Milne, \emph{Jacobian Varieties}, Chapter VII of
\emph{Arithmetic Geometry}, Springer, 1986; author's corrected version
of 12 June 2021, Theorem 6.6 and Summary 6.11. \url{https://www.jmilne.org/math/xnotes/JVs.pdf}.

\bibitem{Bost}
J.-B.~Bost, \emph{Algebraic leaves of algebraic foliations over number
fields}, Publ. Math. IHES 93 (2001), 161--221,
\S4.1.4, Props.~4.5--4.6, equations (4.11), (4.14).
\url{https://www.numdam.org/article/PMIHES_2001__93__161_0.pdf}.

\bibitem{Fulton}
W.~Fulton, \emph{Intersection Theory}, second edition, Springer, 1998,
Th.~3.3(b), projective bundle formula.

\bibitem{GilletSoule}
H.~Gillet and C.~Soul\'e, \emph{Arithmetic intersection theory},
Publ. Math. IHES 72 (1990), 93--174, \S4.2.11, \S4.3,
Th.~4.3.2.
\url{https://www.numdam.org/article/PMIHES_1990__72__93_0.pdf}.

\bibitem{Moriwaki}
A.~Moriwaki, \emph{Hodge index theorem for arithmetic cycles of
codimension one}, Theorems A and B, arXiv:alg-geom/9403011v4.
\url{https://arxiv.org/pdf/alg-geom/9403011}.

\bibitem{Kunnemann}
K.~K\"unnemann, \emph{Some remarks on the arithmetic Hodge index
conjecture}, Compositio Math. 99 (1995), 109--128,
Th.~4.1 and \S5.
\url{https://www.numdam.org/article/CM_1995__99_2_109_0.pdf}.

\bibitem{Shioda}
T.~Shioda, \emph{On the Mordell--Weil lattices}, Comment. Math. Univ.
St. Pauli 39 (1990), 211--240, Lemma 8.1, Ths.~8.4, 8.6,
equations (8.8), (8.11)--(8.14).
\url{https://rikkyo.repo.nii.ac.jp/record/10027/files/AA00610867_39-02_05.pdf}.

\bibitem{Wilson}
R.~M.~Wilson, \emph{A diagonal form for the incidence matrices of
$t$-subsets vs.\ $k$-subsets}, European J. Combin. 11 (1990),
609--615. Used through the verified Corollary 3 of the next reference.

\bibitem{Ghorbani}
E.~Ghorbani, G.~B.~Khosrovshahi, Ch.~Maysoori, M.~Mohammad-Noori,
\emph{Inclusion Matrices and Chains}, arXiv:0709.3144v1,
Corollary 3, p.~10; J. Combin. Theory Ser. A 115 (2008), 878--887.
\url{https://arxiv.org/pdf/0709.3144}.

\bibitem{Mok}
C.~P.~Mok, \emph{Endoscopic Classification of Representations of
Quasi-Split Unitary Groups}, arXiv:1206.0882v5,
\S2.2, equations (2.2.9)--(2.2.10), and \S2.3, equation (2.3.1).
\url{https://arxiv.org/pdf/1206.0882}.

\bibitem{MilneEtale}
J.~S.~Milne, \emph{Lectures on \'Etale Cohomology}, version 2.21,
Th.~23.3, \S27 and Prop.~27.3.
\url{https://www.jmilne.org/math/CourseNotes/LEC.pdf}.
\end{thebibliography}
\end{document}
